%% RhoForge GPU Compiler Specification
%%
%% Build:
%%   pdflatex rhoforge-gpu-compiler-spec.tex
%%   pdflatex rhoforge-gpu-compiler-spec.tex

\documentclass[11pt]{article}
\usepackage[T1]{fontenc}
\usepackage[utf8]{inputenc}
\usepackage{lmodern}
\usepackage[margin=1in]{geometry}
\usepackage{amsmath,amssymb,amsthm}
\usepackage{booktabs}
\usepackage{array}
\usepackage{tabularx}
\usepackage{xcolor}
\usepackage{microtype}
\usepackage{xurl}
\usepackage[colorlinks=true,linkcolor=blue!55!black,
  citecolor=blue!55!black,urlcolor=blue!55!black]{hyperref}
\setlength{\emergencystretch}{3em}

\newtheorem{requirement}{Requirement}[section]
\newtheorem{invariant}[requirement]{Invariant}
\newtheorem{acceptance}[requirement]{Acceptance gate}

\newcommand{\CeTTa}{\textsc{CeTTa}}
\newcommand{\RhoForge}{\textsc{RhoForge}}
\newcommand{\COMM}{\textsc{comm}}
\newcommand{\code}[1]{\texttt{#1}}
\newcommand{\codepath}[1]{\nolinkurl{#1}}
\newcommand{\NameId}{\ensuremath{\mathsf{NameId}}}
\newcommand{\NodeId}{\ensuremath{\mathsf{NodeId}}}
\newcommand{\CompId}{\ensuremath{\mathsf{CompId}}}

\title{\RhoForge:\\
A Semantics-Preserving GPU Reaction Compiler for \CeTTa}
\author{Oru\v{z}i \and Zarathustra Goertzel}
\date{August 2026 --- design specification}

\begin{document}
\maketitle

\begin{abstract}
This document specifies a GPU backend for the reflective rho-calculus engine
in \CeTTa.  The objective is not to place the existing pointer-based
interpreter on a device.  It is to compile canonical rho processes into a
flat reaction representation in which channel matching is a segmented join,
endpoint ownership is an occurrence-indexed claim, communication is a
verified rewrite, and residual construction is parallel compaction.  Three
execution contracts are kept distinct: one-path scheduled execution, exact
successor-frontier enumeration, and serializable waves of independent
communications.  A static-graph specialization maps event-driven neural
networks and conserved causal-credit learning to fixed-width device records,
allowing inference and local learning without an autograd graph or a PyTorch
runtime.  Correctness is fail-closed: every optimization must agree with the
canonical CPU frontier or emit a replayable reduction certificate.  The
design therefore treats the rho semantics as the specification and the GPU
layout as a refinement, not as a replacement language.
\end{abstract}

\paragraph{Status.}
This is an implementation specification, not a performance report.  It is
grounded in the current \CeTTa{} rho engine and its test suites, but no GPU
backend is claimed to exist yet.  ``Must'', ``shall'', and acceptance gates
below are normative.

\section{Objective and boundary}

The target is a direct path
\[
  \text{rho/\CeTTa{} source}
  \longrightarrow \text{canonical reaction IR}
  \longrightarrow \text{GPU execution}
  \longrightarrow \text{rho residual or frontier}.
\]
For neural workloads, the same path shall carry locally generated activity,
eligibility receipts, success signals, and conserved credit.  There is no
tensor autograd graph, backward tape, optimizer object, or Python runtime in
the required execution path.

The backend is not required to accelerate every rho process.  A single small,
highly reflective process may be faster on the CPU.  The GPU target is the
regime with one or more of the following:
\begin{itemize}
\item many endpoints or many independent process instances;
\item stable topology with changing messages;
\item wide channel-local reaction frontiers;
\item repeated event-driven neural episodes; or
\item cost-accounted waves with substantial independent work.
\end{itemize}

\begin{requirement}[Semantic authority]
The existing strict-core and cost-profile semantics remain authoritative.
No kernel may invent a reaction, omit an enabled reaction in an exact mode,
or identify two occurrences merely because their payloads are equal.
\end{requirement}

\section{Verified baseline in the current engine}

The current implementation already contains the conceptual prototype of a
reaction backend.

\begin{itemize}
\item \codepath{src/atom.h} represents symbols, variables, grounded values,
  and expressions as arena-allocated \code{Atom} nodes.  Expression children
  are pointer arrays; structural hashes and hash-consing are available.
\item \codepath{src/rhocalc_core.c} flattens parallel components, assigns
  alpha-invariant canonical keys, extracts send and receive endpoints, sorts
  endpoints by channel key, and applies \COMM{} substitution.
\item The sequential \code{RhoMachine} rebuilds its component and endpoint
  index from the residual process each round.
\item The threaded strict-core path publishes endpoints into per-channel
  buckets and produces a residual corresponding to some legal reduction
  path.
\item The cost profile constructs compatible waves and atomically claims
  endpoint and funding occurrences before commit.  Causal receipts observe
  the same resource consumption rather than changing it.
\item Canonical and rotating schedulers, exact successor frontiers, bounded
  execution status, differential tests, and rho-specific benchmarks already
  provide reference implementations for a device backend.
\end{itemize}

The GPU compiler should preserve this semantic factoring while replacing
heap strings, pointer chasing, and whole-residual reindexing on its hot path.

\section{Canonical device reaction IR}

\subsection{Host canonicalization}

Parsing, profile validation, symbol interning, and the first canonicalization
remain host responsibilities in the initial backend.  The host lowers bound
variables to a stable locally nameless form, flattens top-level parallel
composition, orients the name equation $@(*x)=x$, and interns every canonical
quoted name.

\begin{invariant}[Collision-safe identity]
A \NameId{} is never accepted from a hash alone.  The host maintains a
canonical byte representation and resolves collisions by equality checking.
The device may use a compact 64- or 128-bit identifier only after this table
has certified the mapping.
\end{invariant}

\subsection{Immutable term DAG}

The initial ABI uses immutable nodes and contiguous child storage.

\begin{center}
\begin{tabularx}{\textwidth}{@{}l l X@{}}
\toprule
Buffer & Element type & Meaning\\
\midrule
\code{nodeTag} & \code{u8} & nil, par, send, receive, quote, drop, value\\
\code{nodeAux} & \code{u64} & symbol, variable, value, or profile payload\\
\code{childOffset} & \code{u32/u64} & first child in contiguous child buffer\\
\code{childCount} & \code{u32} & arity\\
\code{children} & \NodeId{} & immutable DAG edges\\
\code{nameOfQuote} & \NameId{} & certified canonical name for quoted nodes\\
\bottomrule
\end{tabularx}
\end{center}

Widths are selected at compile time from a checked capacity header.  A graph
that exceeds the selected width is rejected or recompiled with the wider ABI;
wraparound is forbidden.

\subsection{Occurrence layer}

Terms and occurrences are different objects.  Equal terms may occur multiple
times and must remain separately consumable.

\begin{center}
\begin{tabularx}{\textwidth}{@{}l l X@{}}
\toprule
Field & Type & Meaning\\
\midrule
\code{componentId} & \CompId{} & unique top-level occurrence\\
\code{componentRoot} & \NodeId{} & immutable term DAG root\\
\code{generation} & \code{u32/u64} & stale-reference protection\\
\code{alive} & bitset & occurrence is present in current residual\\
\code{endpointKind} & bit & send or receive\\
\code{channelId} & \NameId{} & canonical channel identity\\
\code{bodyOrPayload} & \NodeId{} & send payload or receive continuation\\
\code{binderMeta} & packed word & substitution depth and profile metadata\\
\bottomrule
\end{tabularx}
\end{center}

The endpoint table is structure-of-arrays.  A radix sort orders it first by
channel, then by endpoint kind and component identifier; alternatively, a
persistent channel index populates it directly.  Segment offsets give a
CSR-like channel directory.

\begin{invariant}[Multiplicity]
Every live send and receive occurrence has exactly one endpoint record, and
every endpoint record refers to exactly one live component generation.
\end{invariant}

\section{Kernel pipeline}

The reference device round consists of six stages.

\paragraph{K0: dirty-component classification.}
Classify newly produced components, emit endpoint records, and identify
non-endpoint residuals.  The first implementation may classify the entire
residual; the production path updates only dirty generations.

\paragraph{K1: channel grouping.}
Sort endpoint records or insert them into a persistent segmented channel
table.  Produce send and receive ranges for every live channel.

\paragraph{K2: reaction planning.}
Apply the selected execution contract.  A plan names exact send and receive
occurrences, the channel, continuation, payload, and any funding occurrences.
Planning does not mutate the residual.

\paragraph{K3: occurrence claim.}
Atomically claim every occurrence in a plan.  Claims use component identity
and generation, not term equality.  Multi-resource plans acquire claims in a
canonical order.  A partial claim is rolled back; it is never exposed as a
committed reaction.

\paragraph{K4: substitution and emission.}
Perform capture-avoiding substitution in the receive body and normalize the
result according to the chosen rho profile.  Immutable subterms are shared.
New nodes are allocated through prefix-summed per-plan sizes or a checked
device arena.  Allocation failure is an explicit status.

\paragraph{K5: commit and compaction.}
Consume claimed occurrences, append generated components, compact dead
records, update dirty channel segments, and optionally emit a reduction
certificate.  The result is a valid residual even if execution stops here.

\begin{acceptance}[Single-round parity]
For every strict-core fixture in the admitted differential corpus, decoding
the device successor set shall equal the canonical CPU successor set as a set
of canonical residuals.  Multiplicity-sensitive tests shall additionally
agree on occurrence consumption.
\end{acceptance}

\section{Execution contracts must remain separate}

\subsection{Canonical one-path mode}

This mode selects exactly the reaction chosen by the canonical CPU scheduler.
It is the primary debugging oracle and may intentionally expose limited GPU
parallelism.  Parallel work occurs within indexing, substitution, and
residual construction, not by changing the selected path.

\subsection{Rotating one-path mode}

The rotating scheduler preserves the current turn state and selection rule.
The device result must serialize to the same scheduled path.  A different
legal rho path is not sufficient for scheduler parity.

\subsection{Exact frontier mode}

All distinct enabled pairings are enumerated and canonical duplicate
residuals are collapsed exactly as in the reference engine.  Search-budget or
memory exhaustion is returned explicitly; it must not be reported as
quiescence or as a complete frontier.

\subsection{Serializable wave mode}

A wave may commit several reactions only when their claimed occurrence sets
are disjoint.  If alternative pairing can change an observable residual, the
wave planner must either enumerate the alternatives or fall back to a one-path
contract.  Every committed wave emits an ordering witness or passes a
serializability check showing that it corresponds to a legal sequence of
ordinary \COMM{} steps.

\begin{requirement}[No semantic performance flag]
Thread count, block size, and device selection may change performance but
shall not silently switch among canonical, rotating, frontier, and wave
semantics.
\end{requirement}

\section{Cost-profile extension}

The cost profile extends each reaction plan with its exact funding cover.
Purse and signature occurrences enter the same occurrence-claim protocol as
the communication endpoints.

\begin{invariant}[Atomic funded commit]
A cost reaction commits iff it owns the selected receive, selected send, and
every funding occurrence in its certified cover.  No residual exposes a state
in which endpoints were consumed but funding was only partly acquired.
\end{invariant}

Funding-cover selection, signature multiplication, event identifiers,
producer causes, and ordered receipt material must agree with the CPU cost
profile.  State-only and receipt-observed device executions shall consume the
same occurrences.

The first GPU milestone may exclude the cost profile, but its IR must preserve
component identity and generation from the start; otherwise atomic funding
cannot be added without redesign.

\section{Static reaction graphs for neural computation}

\subsection{Compilation model}

An event-driven neural network is compiled as a mostly static reaction graph:
channels and handlers are immutable; spikes, traces, signals, and credit
ownership are dynamic messages.  A batch dimension represents independent
networks, episodes, particles, or environments.  Batching independent graphs
is the principal source of occupancy when one graph has a sparse frontier.

The compiler emits:
\begin{itemize}
\item a stable table of neuron, synapse, and resource channel IDs;
\item immutable handler continuations;
\item per-instance dynamic occurrence buffers;
\item optional fixed-capacity queues for bounded-state models; and
\item observation kernels deriving outputs and weights from messages.
\end{itemize}

\subsection{Conserved causal-credit fast path}

The conserved-credit learner has a particularly regular image.  A receipt
slot, signal slot, and credit slot each occupy one finite channel state.  Unary
actions are one-token replacements.  Redemption consumes three state tokens
and publishes their three successors through three ordinary \COMM{} steps.
Derived synaptic efficacy is
\[
  w_k=b+q\,\#\{i\mid\operatorname{owner}(i)=k\},
\]
so no dense gradient tensor or optimizer state is required.

The general backend shall execute the literal compiled rho protocol.  An
optional fused redemption kernel may replace the three administrative steps
only after a refinement proof establishes the same residual, resource
consumption, failure behavior, and observable trace under the admitted
scheduler.

\begin{invariant}[No hidden learning channel]
The device fast path may cache counts or derived weights, but the authoritative
learning state remains the finite receipt, signal, and credit-resource image.
A cache discrepancy is a correctness failure, not an alternative optimizer.
\end{invariant}

\subsection{What bypassing PyTorch means}

The required runtime consists of a native host compiler, a CUDA/HIP-style
device backend, and a small C ABI.  Python bindings may be added for
experimentation but are not on the required execution path.  There is no
autograd, optimizer step, tensor module graph, or framework allocator in the
steady-state loop.  Dense numerical kernels may still be called where a
process payload explicitly denotes a matrix operation; that is an ordinary
grounded operation, not the definition of learning.

\section{Correctness and certification}

Let $\mathcal{C}$ compile a well-formed rho process to device state and let
$\mathcal{D}$ decode a device residual to a canonical rho process.  The proof
and testing programme has five layers.

\begin{enumerate}
\item \textbf{Representation.}
  $\mathcal{D}(\mathcal{C}(P))=\operatorname{canon}(P)$, including occurrence
  multiplicity and quoted-name identity.
\item \textbf{Single-plan soundness.}
  Every committed device plan decodes to one legal rho \COMM{} step, or to the
  declared finite administrative sequence of a proved fused primitive.
\item \textbf{Frontier completeness.}
  Exact mode contains every CPU successor and no additional successor.
\item \textbf{Wave serializability.}
  Every wave decodes to some sequence of legal single steps with the same
  residual and receipt observation.
\item \textbf{Run status.}
  Quiescent, reduction-limit exhausted, search-budget exhausted, allocation
  failure, and validation failure remain distinct results.
\end{enumerate}

A compact reaction certificate records profile, scheduler contract, input
residual digest, selected occurrence IDs and generations, channel IDs,
funding IDs, emitted roots, and output residual digest.  A CPU checker replays
the certificate against canonical semantics.  The Lean refinement target is
a decoder theorem connecting checked certificates to the established
\code{Reduces} relation.

\begin{acceptance}[Fail-closed differential gate]
No optimized kernel is admitted by throughput alone.  It must pass canonical
single-step differential tests, scheduler-specific run tests, multiplicity
tests, malformed-state rejection, exhaustion-status tests, and randomized
small-process frontier comparison.
\end{acceptance}

\section{Runtime API}

The minimal native interface is intentionally small.

\begin{verbatim}
rf_context_create(device, options, *ctx)
rf_compile_rho(ctx, canonical_bytes, profile, *program)
rf_instantiate(program, batch_size, initial_messages, *state)
rf_step(state, execution_contract, budget, *result)
rf_run(state, execution_contract, limits, *result)
rf_decode_residual(result, instance, *canonical_bytes)
rf_export_certificate(result, *bytes)
rf_observe(result, observer_id, *buffer)
rf_destroy(handle)
\end{verbatim}

Every call returns a typed status.  Program objects are immutable and
shareable; state objects own dynamic occurrences.  Host/device transfer is
explicit.  Observers cannot mutate the reaction state.

\section{Milestones and acceptance gates}

\begin{center}
\begin{tabularx}{\textwidth}{@{}l X X@{}}
\toprule
Milestone & Deliverable & Admission gate\\
\midrule
M0 & Versioned flat IR and CPU decoder & Canonical round trip; multiplicity\\
M1 & GPU endpoint extraction and channel join & Exact endpoint parity\\
M2 & Strict-core canonical one-step backend & Successor differential corpus\\
M3 & Rotating and exact-frontier modes & Scheduler/frontier parity; exhaustion\\
M4 & Incremental dirty-channel index and compaction & Same gates plus speed/memory evidence\\
M5 & Cost-profile atomic claims and receipts & Cost differential and causal replay\\
M6 & Static neural graph compiler & Event-trace parity with CPU rho\\
M7 & Conserved-credit device fast path & Literal-vs-fused refinement and task parity\\
M8 & Lean certificate bridge & Checked certificate implies rho reduction\\
\bottomrule
\end{tabularx}
\end{center}

Performance gates begin only after M2 correctness.  Required measurements
include reactions per second, bytes moved per committed reaction, endpoint
index cost, substitution cost, residual-compaction cost, active-lane fraction,
batch scaling, and host/device transfer.  Comparisons include the current
sequential and threaded rho engines and the existing contention, fanout,
pipeline, and cost-parallel workloads.  A PyTorch comparison is meaningful
only for a matched neural task and excludes import/startup time unless both
systems are evaluated end to end.

\section{Risks and required fallbacks}

\paragraph{Irregular reflection.}
Dynamic quotation and process generation may produce divergent graph growth.
The backend must support checked device allocation and a semantic CPU fallback
at a round boundary.  Falling back is preferable to truncating a process.

\paragraph{Small frontier under-utilization.}
A single sparse graph may not occupy a GPU.  Batch independent graphs, keep a
persistent kernel for repeated episodes, or use the CPU path.  The compiler
may choose a backend from static and observed frontier width, provided the
semantic contract is unchanged.

\paragraph{Canonicalization cost.}
String canonical keys are unsuitable as the device hot representation.
Verified name interning, structural hashes with collision checks, and
incremental dirty-channel maintenance are required.  Host canonical bytes
remain the audit representation.

\paragraph{Nondeterministic atomics.}
Atomic arrival order must not define canonical or rotating semantics.  Those
modes plan deterministically before claim.  A throughput wave may use atomics
only within its explicitly weaker serializable-path contract.

\paragraph{Fused-kernel semantic drift.}
No neural or cost fast path is accepted because it is ``obviously the same.''
Literal rho execution is the reference, and the fused path requires residual,
failure, resource, and observation equivalence.

\section{Non-goals}

The first backend does not attempt to compile arbitrary grounded C functions
to device code, guarantee speedup for every process, replace the host parser,
or treat floating-point nondeterminism as rho nondeterminism.  It does not
define a new neural-network formalism.  It compiles a subset of already
declared grounded operations and the rho communication substrate, with
explicit rejection or host fallback outside that subset.

\section{Conclusion}

The formatting is the optimization.  Once reflective terms are separated
from occurrences, canonical names from heap strings, and channel matching
from recursive traversal, rho execution becomes a sequence of familiar data-
parallel operations: classify, group, plan, claim, rewrite, and compact.  The
same representation makes a process-native neural engine plausible.  Static
handlers describe the graph; messages describe activity and causal learning;
and conserved credit replaces an optimizer state with explicit finite
resources.

The design remains rho rather than merely rho-inspired because every device
transition is checked against the existing semantics.  That constraint is
also the engineering advantage: the CPU engine, differential corpus, causal
receipts, and Lean reduction relation together provide unusually strong
oracles for building an aggressive GPU backend without guessing what the
program meant.

\end{document}
